\subsubsection*{\texttt{31-retrieval-fusion-2q}}
\label{app:samples:31_retrieval_fusion_2q}

\paragraph{Q3: How long does the Haste spell last?}
\noindent \textbf{Corr.:}~4.00/5\quad \textbf{Faith.:}~1.00\quad \textbf{Relev.:}~1.00\quad \textbf{Ctx.Rel.:}~1.00\\

\noindent\textbf{Antwort:} With a 3rd-level spell slot, the Haste spell lasts for concentration, up to 1 minute.

\noindent\textbf{Abgerufene Chunks:}
\begin{itemize}[noitemsep, topsep=2pt]
  \item[\textbf{\#1}] \texttt{[0.0333]} \textit{Spells/Haste} -- \#\#\# Haste  *3rd-level transmutation*  **Casting Time:** 1 action  **Range:** 30 \ldots
  \item[\textbf{\#2}] \texttt{[0.0328]} \textit{Spells/# Spellcasting} -- \#\# Duration  A spell's duration is the length of time the spell persists. A dura\ldots
  \item[\textbf{\#3}] \texttt{[0.0317]} \textit{Spells/Geas} -- \#\#\# Geas  *5th-level enchantment*  **Casting Time:** 1 minute  **Range:** 60 fee\ldots
  \item[\textbf{\#4}] \texttt{[0.0313]} \textit{Classes/Sorcerer} -- \#\#\#\# Extended Spell  When you cast a spell that has a duration of 1 minute or lo\ldots
  \item[\textbf{\#5}] \texttt{[0.0311]} \textit{Spells/Modify Memory} -- \#\#\# Modify Memory  *5th-level enchantment*  **Casting Time:** 1 action  **Range:\ldots
\end{itemize}
\vspace{4pt}
\hrule
\vspace{6pt}

\paragraph{Q5: Can a grappled creature cast spells?}
\noindent \textbf{Corr.:}~1.00/5\quad \textbf{Faith.:}~0.00\quad \textbf{Relev.:}~0.00\quad \textbf{Ctx.Rel.:}~0.25\\

\noindent\textbf{Antwort:} The context does not provide information on whether a grappled creature can cast spells or not. The first context only mentions that a grappled creature's speed becomes 0 and it can't benefit from any bonus to its speed. The second context is about spellcasting for monsters, but it doesn't mention a\ldots

\noindent\textbf{Abgerufene Chunks:}
\begin{itemize}[noitemsep, topsep=2pt]
  \item[\textbf{\#1}] \texttt{[0.0323]} \textit{Gamemastering/Conditions} -- \#\# Grappled  - A grappled creature's speed becomes 0, and it can't benefit from \ldots
  \item[\textbf{\#2}] \texttt{[0.0323]} \textit{Gameplay/Combat} -- \#\#\# Grappling  When you want to grab a creature or wrestle with it, you can use \ldots
  \item[\textbf{\#3}] \texttt{[0.0323]} \textit{Classes/Sorcerer} -- \#\#\#\# Careful Spell  When you cast a spell that forces other creatures to make a \ldots
  \item[\textbf{\#4}] \texttt{[0.0318]} \textit{Gameplay/Combat} -- \#\# Cast a Spell  Spellcasters such as wizards and clerics, as well as many monst\ldots
  \item[\textbf{\#5}] \texttt{[0.0313]} \textit{Spells/Arcane Hand} -- ***Grasping Hand***. The hand attempts to grapple a Huge or smaller creature wit\ldots
\end{itemize}
\vspace{4pt}
\hrule
\vspace{6pt}

\paragraph{Q12: If a creature has exhaustion level 3 and is also restrained, what is its movement speed?}
\noindent \textbf{Corr.:}~5.00/5\quad \textbf{Faith.:}~1.00\quad \textbf{Relev.:}~1.00\quad \textbf{Ctx.Rel.:}~1.00\\

\noindent\textbf{Antwort:} The context provided does not give enough information to determine the movement speed of a creature with both the exhaustion and restrained conditions. The rules for the restrained condition state that the creature's speed becomes 0, but the exhaustion condition only halves the creature's speed. How\ldots

\noindent\textbf{Abgerufene Chunks:}
\begin{itemize}[noitemsep, topsep=2pt]
  \item[\textbf{\#1}] \texttt{[0.0333]} \textit{Gamemastering/Conditions} -- \#\# Restrained  - A restrained creature's speed becomes 0, and it can't benefit f\ldots
  \item[\textbf{\#2}] \texttt{[0.0318]} \textit{Gamemastering/Conditions} -- \#\# Exhaustion  Some special abilities and environmental hazards, such as starvat\ldots
  \item[\textbf{\#3}] \texttt{[0.0301]} \textit{Gamemastering/Conditions} -- \#\# Unconscious  - An unconscious creature is incapacitated (see the condition), \ldots
  \item[\textbf{\#4}] \texttt{[0.0301]} \textit{Monsters/Lion (Creature)} -- \#\#\# Lion  *Large beast, unaligned*  **Armor Class** 12  **Hit Points** 26 (4d10+\ldots
  \item[\textbf{\#5}] \texttt{[0.0164]} \textit{Gameplay/Adventuring} -- \#\# Suffocating  A creature can hold its breath for a number of minutes equal to \ldots
\end{itemize}
\vspace{4pt}
\hrule
\vspace{6pt}

\paragraph{Q16: Can you use Action Surge to cast two leveled spells in one turn?}
\noindent \textbf{Corr.:}~2.00/5\quad \textbf{Faith.:}~0.00\quad \textbf{Relev.:}~0.00\quad \textbf{Ctx.Rel.:}~0.25\\

\noindent\textbf{Antwort:} I cannot find this in the provided rules that you can use Action Surge to cast two leveled spells in one turn. Action Surge only allows you to take an additional action on top of your regular action and a possible bonus action. Casting a spell is considered an action, so you can only cast one spell \ldots

\noindent\textbf{Abgerufene Chunks:}
\begin{itemize}[noitemsep, topsep=2pt]
  \item[\textbf{\#1}] \texttt{[0.0331]} \textit{Classes/Fighter} -- \#\#\# Action Surge  Starting at 2nd level, you can push yourself beyond your norma\ldots
  \item[\textbf{\#2}] \texttt{[0.0310]} \textit{Characterizations/Multiclassing} -- If you have more than one spellcasting class, this table might give you spell sl\ldots
  \item[\textbf{\#3}] \texttt{[0.0306]} \textit{Classes/Wizard} -- \#\#\#\# Learning Spells of 1st Level and Higher  Each time you gain a wizard level,\ldots
  \item[\textbf{\#4}] \texttt{[0.0296]} \textit{Classes/Paladin} -- \#\#\# Extra Attack  Beginning at 5th level, you can attack twice, instead of once,\ldots
  \item[\textbf{\#5}] \texttt{[0.0167]} \textit{Spells/Confusion} -- \#\#\# Confusion  *4th-level enchantment*  **Casting Time:** 1 action  **Range:** 9\ldots
\end{itemize}
\vspace{4pt}
\hrule
\vspace{6pt}
